% ============================================================
% CHAPTER 8: CONCLUSION AND FUTURE WORK
% ============================================================
\chapter{Conclusion and Future Work}

This chapter synthesises the research contributions, provides formal answers to the three guiding research questions, discusses practical implications for the UAV security community, and outlines directions for future investigation.

\section{Summary of Contributions}

This research has demonstrated that software-only GPS spoofing detection for PX4-based unmanned aerial vehicles is feasible using cross-sensor consistency checks and machine learning. The specific contributions to knowledge and practice are:

\begin{enumerate}
    \item \textbf{A physics-based GPS--IMU consistency discriminator} that reliably identifies GPS spoofing by detecting the fundamental physical contradiction between GPS-reported position change and IMU-observed motion. The discriminator is independent of attack waveform characteristics, operates on standard telemetry without hardware augmentation, and provides a principled, model-free detection signal.
    \item \textbf{Eight heuristic labelling rules} that enable automated training data generation from unlabelled flight logs, addressing the annotation bottleneck that has limited the deployment of supervised ML approaches for GPS security. The heuristic framework is extensible: new detection heuristics can be added as additional sensor modalities or attack types are discovered.
    \item \textbf{A complete end-to-end detection pipeline} encompassing: multi-threaded MAVLink telemetry acquisition at 10\,Hz with thread-safe state management; five-stage data preprocessing (cleaning, feature engineering, labelling, windowing, splitting); dual-model training with Random Forest (100\% test accuracy) and 1D CNN (83.3\% test accuracy); streaming inference with EMA smoothing and terrain-aware model dispatching; and graduated post-detection response via the GPS Trust Index.
    \item \textbf{Two fully functional user interfaces}: a Textual-based TUI management console providing component lifecycle management with PID tracking, integrated log viewing, and one-click attack injection; and a Streamlit web dashboard with live GPS Trust Index display, anomaly probability monitoring, and historical log replay.
    \item \textbf{A reproducible Docker-based simulation environment} using PX4 SITL and Gazebo Harmonic, enabling controlled experimentation with scripted GPS spoofing attacks under varied flight profiles. The environment, configuration, and all software components are released as open-source resources to support replication and extension by the research community.
    \item \textbf{The ``CNN Paradox'' analysis}---an empirical finding that the 1D CNN achieves perfect recall but 100\% false positive rate on the limited normal sample---which provides a methodological caution for ML-based GPS security research: classification metrics must be evaluated with explicit analysis of confusion structure, and class-imbalanced datasets demand stratified evaluation protocols.
\end{enumerate}

\section{Answers to Research Questions}

\textbf{RQ1: What are the primary vulnerabilities of GPS in autonomous mobile robots to spoofing attacks?}

The primary vulnerability of GPS in autonomous mobile robots---specifically in the PX4 flight stack---is the EKF2's architectural dependence on GPS as a high-weight, high-confidence observation channel. The EKF2 fuses GPS position updates weighted by their self-reported accuracy metrics (eph, epv). Because spoofed signals deliberately report low position error and healthy fix quality, the filter accepts falsified data without immediate rejection. The natural protection mechanism---innovation gating, where observations exceeding expected residuals are rejected---only triggers when the discrepancy between predicted and observed position has grown sufficiently large, which may require tens of seconds depending on the specific innovation gate configuration. During this period, the drone's control loops are acting on progressively more corrupted position estimates. Countermeasures that rely on signal authentication or hardware augmentation are unavailable on the vast installed base of civilian receivers, creating a systemic vulnerability that can only be addressed through software-based cross-validation with other sensor modalities.

\textbf{RQ2: How can sensor fusion techniques effectively detect GPS spoofing in real time?}

Sensor fusion---specifically, the systematic cross-validation of GPS-reported position against IMU-measured physical motion---provides an effective real-time spoofing detection capability. The GPS--IMU consistency score (\(\Delta \Phi\)), defined as the Euclidean discrepancy between IMU-integrated displacement and GPS-reported displacement over a time window, constitutes a discriminative signal that distinguishes spoofing (where GPS position changes without corresponding IMU-detected motion) from natural flight disturbances such as wind gusts (where both GPS and IMU register the physical displacement). The heuristic implementation of this principle, which flags conditions where IMU angular rates remain below 0.05\,rad/s while GPS reports speed exceeding 5\,m/s, runs with sub-millisecond per-sample latency and requires no prior training data. When combined with Random Forest classification on 30-sample sliding windows of full PX4 telemetry (34 features), the system achieves 100\% detection accuracy on the evaluated test set. The feasibility of real-time detection is therefore affirmed, with the caveat that the 3-second window accumulation delay exceeds theoretical safe-reaction times for high-speed obstacle-avoidance scenarios, motivating future work on reduced-latency architectures as discussed below.

\textbf{RQ3: How can a dynamic GPS Trust Index be designed to quantify the reliability of GPS data?}

The GPS Trust Index (GTI) is designed as a temporally smoothed, sigmoid-transformed function of the GPS--IMU consistency score \(\Delta \Phi_t\) (Equation~\ref{eq:gti_full}). The design encodes three key properties: (a) \textit{continuity}---the index is a real number in \([0, 1]\), enabling graduated rather than binary trust decisions; (b) \textit{temporal smoothing}---an exponential moving average with \(\alpha = 0.95\) prevents transient disturbances from causing abrupt trust degradation while allowing persistent anomalies to drive the index toward zero; and (c) \textit{probabilistic interpretability}---the sigmoid maps the unbounded physical inconsistency \(\Delta \Phi\) to a trust value interpretable as the estimated probability that the current GPS position update is authentic. The GTI enables a graduated response protocol that escalates protective measures (sensor weight modulation, hold mode activation, emergency return-to-launch) as trust declines, balancing operational safety with mission continuity. The design is modular in that additional trust-influencing factors (e.g., satellite geometry, innovation monitoring, consistency with magnetometer heading) can be incorporated into the sigmoid argument as additive terms without changing the underlying architecture.

\section{Practical Implications}

This research carries several practical implications for the UAV security community. First, it demonstrates that effective GPS spoofing detection is achievable on existing PX4 hardware without firmware modifications---a departure from prior work that assumed the necessity of hardware upgrades or access to encrypted military signals. The companion-computer architecture (a Raspberry Pi or similar single-board computer running the detection software alongside the flight controller) is compatible with current commercial UAV configurations.

Second, the modular pipeline design allows operators to select appropriate sensitivity and response parameters for their specific operational context. A drone surveying agricultural fields at low altitude and low speed can tolerate longer detection latency and higher false positive rates than a drone inspecting power lines at close proximity. The GTI threshold and EMA parameters can be tuned per mission profile.

Third, the automated labelling pipeline enables continuous model improvement as new flight data becomes available. Each operational flight generates unlabelled telemetry that can be processed by the heuristic labeller to produce training data, enabling the ML models to adapt to platform-specific sensor characteristics and operational patterns without manual annotation effort.

Fourth, the open-source availability of all system components---simulation environment, data capture, preprocessing, training, inference, and interfaces---facilitates adoption, extension, and independent verification by other researchers and practitioners.

\section{Future Work}

Several directions for future research are motivated by the limitations and open questions identified in this study:

\begin{enumerate}
    \item \textbf{Physical hardware validation}: The most critical next step is to deploy and evaluate the detection system on physical PX4 hardware (e.g., a Pixhawk 4 or Cube Orange flight controller on an actual quadcopter airframe). This would expose the detector to genuine sensor noise characteristics, real RF propagation and multipath effects, and the diverse failure modes that arise in field conditions. Controlled outdoor experiments with a software-defined radio-based GPS spoofer would provide the most rigorous validation of the detection approach.
    \item \textbf{Breaking point analysis}: A systematic parametric sweep over spoofing attack parameters---drift rate (0.1 to 10\,m/s), position offset magnitude (1 to 1000\,m), onset duration (instant to 60-second ramp), and intermittent attack duty cycle---is needed to characterise the detection performance surface. The \textit{breaking point}, defined as the minimum spoofing intensity at which detection probability falls below 90\%, would establish the operational envelope of the system.
    \item \textbf{GTI integration into PX4 EKF2}: The GPS Trust Index should be integrated directly into the open-source PX4 firmware as a modulation factor on the GPS observation innovation covariance matrix. This would enable native, zero-latency GPS trust modulation without requiring a companion computer for detection. The integration would require modifications to the EKF2 observation fusion logic and the addition of the GTI update Equation~\ref{eq:gti_full} to the sensor management framework.
    \item \textbf{Larger balanced datasets}: The current 76-window dataset is modest and has extreme class imbalance. Future datasets should include at least 500 windows with balanced proportions of normal and spoofed flight, multiple distinct flight profiles (hover, cruise, climb, descend, aggressive turns), varied environmental conditions (calm, windy, gusty), and diverse spoofing attack types (drift, meaconing, intermittent, multi-sensor). Such a dataset would enable statistically robust evaluation and model comparison.
    \item \textbf{Reduced-latency architectures}: The 3-second window delay motivates exploration of alternative detection architectures. Single-sample detectors based on instantaneous GPS--IMU consistency (the heuristic H8 provides this already, albeit with higher false positive rates), recurrent neural networks with stateful memory that can detect attacks from shorter observation sequences, attention-based models that focus on the most discriminative temporal positions within a window, and transformer architectures adapted to multivariate sensor streams all warrant investigation for latency-critical applications.
    \item \textbf{Multi-drone cooperative detection}: In swarm or formation-flight operations, multiple drones could share GPS trust information, cross-validate spatial consistency between vehicles, and collectively detect regional spoofing attacks. A drone whose GPS position diverges from the consensus of the swarm would be identified as potentially spoofed. This cooperative detection paradigm could substantially increase the effective detection range and robustness.
    \item \textbf{Cross-platform generalisation}: Extending the detection pipeline to Ardupilot, BetaFlight, and other open-source flight stacks would evaluate the generalisability of the GPS--IMU consistency approach beyond the PX4 ecosystem. The fundamental physical principle is platform-independent, but differences in EKF implementation, sensor fusion weighting, and telemetry message structures would require adaptation.
    \item \textbf{Formal adversarial robustness evaluation}: A rigorous evaluation against \textit{adaptive} spoofing attacks---where the attacker has knowledge of the detection algorithm and designs the spoofing waveform to minimise the GPS--IMU consistency score---is needed to establish the security guarantees of the approach under a threat model that includes detector-aware adversaries.
\end{enumerate}

\section{Concluding Remarks}

GPS spoofing is one of the most insidious threats facing the rapidly expanding unmanned aerial vehicle ecosystem, precisely because it exploits the trust relationship between a critical navigation sensor and the flight control system that depends on it. This research has demonstrated that the same physics that makes GPS spoofing difficult to detect---its ability to produce internally consistent, apparently healthy position estimates---also creates its fundamental detection vulnerability: the inescapable inconsistency between falsified position data and the drone's true physical motion, as measured by its inertial sensors.

The detection system presented in this thesis, combining principled physics-based discrimination with complementary machine learning classification, provides a proof of concept that software-only spoofing protection is achievable on current-generation commercial UAV platforms. As drones become increasingly integrated into civilian airspace, critical infrastructure operations, and autonomous logistics networks, the need for robust, accessible, and deployable spoofing countermeasures will only grow. This work contributes a foundational building block toward that goal, with the explicit intention that its open-source release and systematic evaluation framework will support continued progress by the broader research community.
